\section{Introduction}
\label{sec:intro}

The ultimate goal of robotics should not be to develop specialized skills for specific hardware, but to build a general embodied agent~\cite{pi_0.5,shukor2025smolvla,eo1,zheng26}---similar to the human brain---that can operate across diverse robotic bodies, thereby realizing the vision of ``one mind, multiple forms''. Such an agent must perceive its environment, understand natural language instructions, reason about actions, and learn from interaction~\cite{rt2,liu2026rdt2,cheang2025gr,zheng25}. Despite decades of progress, this vision remains unrealized. Embodied intelligence has yet to reach its transformative moment.
% 机器人领域的最终目标不应是为特定硬件开发专用技能，而应是构建一种通用的具身智能体——类似于人类大脑——能够运行于各种机器人本体之上，从而实现“一脑多形”的愿景。这种智能体必须能感知环境、理解自然语言指令、推理行为并从交互中学习。尽管已有数十年的发展，这一愿景仍未实现，具身智能尚未迎来其突破性时刻。

We identify three key limitations. First, data scale has not reached a critical threshold. Unlike standard vision-language data, embodied data require precise action labels, which are expensive and slow to collect. Most current approaches rely on pre-training sets limited to a single robot type or platform. While effective for narrow deployment, this restricts the development of general policies by limiting exposure to diverse morphologies and task distributions. Second, data quality varies widely and lacks standardization. Action representations, coordinate systems, and control frequencies differ across datasets, making it difficult for models to extract common patterns across embodiments. The resulting fragmentation introduces noise and weakens learning, as models must allocate capacity to memorize idiosyncrasies rather than acquire transferable skills. Third, existing pre-training paradigms are mismatched. Most Vision-Language-Action (VLA) models start from Vision-Language Models (VLMs), whose visual encoders focus on semantic recognition rather than 3D structure or physical dynamics. Yet embodied behavior demands fine spatial understanding and causal reasoning---capabilities that cannot be acquired from 2D supervision alone, nor easily added through mechanisms such as Chain-of-Thought, because such methods operate at the reasoning level without enriching the underlying perceptual representation. Achieving general embodied intelligence therefore requires systematic improvements in data, representation, and training design.
% 我们指出了三个关键限制：第一，数据规模尚未达到临界点。与常规视觉-语言数据不同，具身数据需要精确的动作标注，采集成本高且耗时。现有方法大多依赖单一机器人或平台的数据集进行预训练。虽然适用于局部部署，但这种做法限制了对多样化形态和任务分布的接触，阻碍了通用策略的发展。第二，数据质量参差不齐且缺乏标准化。动作表示、坐标系和控制频率在不同数据集中各不相同，导致模型难以提取跨本体的共性知识。由此产生的碎片化引入了噪声，削弱了学习效果，因为模型不得不将容量用于记忆特异性模式，而非习得可迁移技能。第三，当前的预训练范式存在错配。大多数视觉-语言-动作（VLA）模型基于视觉-语言模型（VLM）初始化，其视觉编码器侧重语义识别，而非三维结构或物理动态建模。然而，具身行为需要精细的空间理解与因果推理能力，这些无法仅通过二维监督获得，也无法通过思维链等机制轻易弥补，因为此类方法作用于推理层面，并未增强底层感知表征。因此，实现通用具身智能需要在数据、表示和训练设计上进行系统性改进。

We address a practical question: given the high cost and hardware dependence of collecting robotic data, can we integrate global open-source datasets---currently isolated in silos---into a unified foundation for training hardware-agnostic VLA models? This suggests a new direction: instead of emerging from closed, proprietary systems, embodied intelligence can grow through the reorganization of public knowledge and incremental capability building, where shared infrastructure enables cumulative progress across research groups.
% 我们探讨一个实际问题：鉴于机器人数据采集的成本高昂且依赖硬件，能否将全球开源数据——目前分散在各个孤岛中——整合为统一基础，用于训练不依赖特定硬件的VLA模型？这指向一种新方向：具身智能不再源于封闭私有系统，而是通过公共知识的重构与能力的逐步积累而成长，其中共享基础设施支持研究团队间的持续进展。

To this end, we introduce \VLAFM{}, a unified approach that jointly optimizes data processing and model architecture. We combine six major open-source datasets---including OXE~\cite{o2024oxe}, OXE-AugE~\cite{ji2025oxeauge}, Agibot-Beta~\cite{agibotbeta}, RoboCoin~\cite{wu2025robocoin}, RoboMind~\cite{wu2024robomind}, and Galaxea~\cite{jiang2025galaxea} --- to form a mixed training set of over 6 million trajectories, the largest such collection outside proprietary efforts. To resolve inconsistencies in action format, coordinate system, and sampling rate, we define a standardized preprocessing pipeline. All actions are converted to delta actions in the end-effector frame; rotations are encoded using rotation vectors to avoid singularities; and we apply a pad-to-dual strategy to support both single-arm and dual-arm tasks within a shared framework. This unified data foundation breaks down barriers between datasets and enables robust cross-embodiment generalization by aligning both geometric and temporal structures across sources.
% 为此，我们提出\VLAFM{}，一种协同优化数据处理与模型架构的统一方法。我们整合六个主要开源数据集——包括OXE、OXE-AugE、Agibot-Beta和RoboCoin——构建了超过600万条轨迹的混合训练集，是目前非私有体系中规模最大的集合。为解决动作格式、坐标系和采样率的不一致，我们定义了标准化预处理流程：所有动作转换为末端执行器坐标系下的增量动作；旋转使用旋转向量编码以避免奇异性；采用pad-to-dual策略，在共享框架下支持单臂与双臂任务。这一统一的数据基础打破了数据壁垒，通过在源间对齐几何与时间结构，实现了强大的跨本体泛化能力。

On the modeling side, we build on the ``VLM + Action Expert'' architecture and introduce two core innovations. First, most VLA models~\cite{GR00T,team2024octo,li2024cogact,reuss2025flower} train the policy to predict noise---either $\epsilon$-prediction or $v$-prediction---which is unstructured, meaningless, and often includes invalid actions such as jittering or discontinuities. These noisy targets spread across high-dimensional space, requiring large model capacity and leading to inefficient learning. Inspired by JiT~\cite{JIT} compilation, we propose the \textit{``Action Manifold Hypothesis''}: successful actions do not scatter randomly but lie on a low-dimensional, smooth manifold shaped by physics, task goals, and environmental constraints. Based on this, we design the Action Manifold Learning (AML) mechanism, where the DiT backbone directly predicts clean action sequences. This shifts the learning objective from fitting noise to projecting onto feasible action manifolds, improving decoding speed and policy stability by reducing output uncertainty. Second, to improve spatial understanding, we add modular components that either inject scene-level structural priors via VGGT~\cite{wang2025vggt} or fuse multi-view observations generated by Qwen-Image-Edit~\cite{wu2025qwenimagetechnicalreport}, allowing flexible enhancement without modifying the main architecture. These modules act as plug-and-play experts that enrich perception with geometry-aware features, addressing the representational gap left by standard VLMs.
% 在模型方面，我们在“VLM + 动作专家”架构基础上提出两项核心创新。首先，多数VLA模型训练策略预测噪声（如ε预测或v预测），这些噪声无结构、无意义，常包含抖动或不连续等无效动作，散布于高维空间，要求模型具备大容量，导致学习效率低下。受JiT启发，我们提出“动作流形假说”：成功动作并非随机分布，而受限于物理规律、任务目标和环境约束，聚集在低维平滑流形上。基于此，我们设计AML机制，使DiT主干直接预测干净的动作序列，将学习目标从拟合噪声转变为向可行流形投影，通过减少输出不确定性，显著提升解码速度与策略稳定性。其次，为增强空间理解，我们引入模块化组件，通过VGGT注入场景级结构先验，或融合LongCat生成的多视角观测特征，可在不修改主干的情况下灵活增强能力。这些模块作为即插即用的专家，通过几何感知特征丰富感知能力，弥补了标准VLM留下的表征差距。

We perform full ablation studies under a consistent training setup, showing that data standardization, architectural changes, and training redesign are orthogonal and their benefits accumulate. Our model achieves average success rates of 98.6\%, 80.5\% and 58.3\% on benchmarks including LIBERO~\cite{liu2023libero}, LIBERO-Plus~\cite{fei25libero-plus}, and RoboCasa GR1 Tabletop Tasks~\cite{GR00T}, outperforming strong baselines such as $\pi_{0.5}$~\cite{pi_0.5}, UniVLA~\cite{bu2025univla}, and openVLA-OFT~\cite{kim2024openvla}. These results confirm a complete pathway---from data to architecture to capability---that enables high-performance, generalizable embodied intelligence through systematic engineering, even without private data. Notably, the gains from each component remain consistent across platforms, suggesting that the learned representations generalize beyond dataset-specific artifacts. We will fully release our data processing and training code to support community-driven progress toward open, shared, and evolving general-purpose robotics.
% 我们在统一训练框架下进行了完整消融实验，结果表明数据标准化、架构改进与训练设计三者相互正交且收益可叠加。模型在Libero、Libero-Plus、和RoboCasa等多个基准上分别取得98.6%、80.5%、58.3%的平均成功率，显著优于pi0、uniVLA和openVLA-OFT等先进方法。结果验证了一条从数据到架构再到能力的完整技术路径，证明即使没有私有数据，也能通过系统工程构建高性能、强泛化的具身智能。值得注意的是，各组件带来的增益在不同平台上保持一致，表明所学表征超越了数据集特异性偏差。我们将全面开源数据处理与训练代码，推动社区共建共享、持续进化的通用机器人生态发展。

